\SourceLeaf{332}

\[
(4), (8), (9) \text{ aus I}
\]
\[
(3), (7), (11) \text{ aus II}
\]
\[
(5), (6), (10) \text{ aus I und II verbunden.}
\]

Jede der Gleichungen (3)—(11) kann zur Auflösung des \textsc{Pothenot}schen Problems dienen, da die zwei Unbekannten, die in jeder Gleichung vorkommen, eine bekannte Summe haben.

Aus Figur II ergeben sich die Viereckgleichungen

\[
\frac{\sin(\mathfrak{A}-A)}{a\alpha} = \frac{\sin(\mathfrak{B}-B)}{b\beta} = \frac{\sin(\mathfrak{C}-C)}{c\gamma} = \frac{rr-\rho\rho}{2\alpha\beta\gamma r}
\]

von selbst, wenn $r$ Halbmesser des Kreises durch $A$, $B$, $C$, $\rho$ Distanz des vierten Punkts vom Mittelpunkte dieses Kreises; und die Gleichheit der Zeichen von

\[
\sin(\mathfrak{A}-A), \qquad \sin(\mathfrak{B}-B), \qquad \sin(\mathfrak{C}-C)
\]

ist das Criterium der physischen Möglichkeit der Aufgabe.

Die geometrische Bedeutung der Formeln (2)—(11) kann sehr zierlich an die fünf Kreise geknüpft werden, die resp.\ durch je drei Punkte von Fig.\ I oder von Fig.\ II beschrieben werden und deren Durchmesser $p$, $q$, $r$, $s$, $t$ sein mögen. Es ist dann

\[
\sin A = \frac{a}{p}, \qquad \sin B = \frac{b}{p}, \qquad \sin C = \frac{c}{p},
\]
\[
\sin\mathfrak{A} = \frac{a}{q}, \qquad \sin 4 = \frac{\gamma}{q}, \qquad \sin 5 = \frac{\beta}{q},
\]
\[
\sin\mathfrak{B} = \frac{b}{r}, \qquad \sin 6 = \frac{\alpha}{r}, \qquad \sin 1 = \frac{\gamma}{r},
\]
\[
\sin\mathfrak{C} = \frac{c}{s}, \qquad \sin 2 = \frac{\beta}{s}, \qquad \sin 3 = \frac{\alpha}{s},
\]
\[
\sin(\mathfrak{A}-A) = \frac{a\alpha}{t}, \qquad \sin(\mathfrak{B}-B) = \frac{b\beta}{t}, \qquad \sin(\mathfrak{C}-C) = \frac{c\gamma}{t}.
\]

Durch Substitution dieser Ausdrücke für die 15 Sinus in den 10 Gleichungen (2)—(11) werden diese zu identischen.

\Needspace{7\baselineskip}Ganz auf dieselbe Weise, wie aus Gleichung (2) die übrigen (3)—(11) identisch folgen, lassen sich auch aus jeder von diesen alle übrigen 9 ableiten. Man braucht z.\,B. nur die Gleichung (3) in die Form zu setzen
